\textbf{1. Kiến thức cần nhớ}
\begin{itemize}
\item \textbf{Quy tắc cộng hai đa thức}: Muốn cộng hai đa thức, ta nối hai đa thức bằng dấu $+$'' rồi bỏ ngoặc và thu gọn các hạng tử đồng dạng.
80493\boxed{A + B}80493
\item \textbf{Quy tắc trừ hai đa thức}: Muốn trừ đa thức $ cho đa thức $, ta lấy $ cộng với đa thức đối của $:
80493\boxed{A - B = A + (-B)}80493
\begin{itemize}
\item Khi bỏ dấu ngoặc có dấu 569X$'' đằng trước, ta phải \textbf{đổi dấu tất cả các hạng tử} trong ngoặc: dấu $+$'' thành 569X$'', dấu 569X$'' thành $+$''.
\end{itemize}
\item \textbf{Tính chất}: Phép cộng đa thức có tính chất giao hoán và kết hợp:
80493A + B = B + A;\quad (A + B) + C = A + (B + C).80493
\end{itemize}

\textbf{2. Phương pháp giải các dạng}

%% DANG 8D1?3-1
\textbf{Dạng 1. Cộng, trừ hai đa thức}
\begin{enumerate}
\item Đặt các đa thức vào trong từng dấu ngoặc đơn:  + (B)$ hoặc  - (B)$.
\item Bỏ dấu ngoặc theo đúng quy tắc đổi dấu.
\item Nhóm các hạng tử đồng dạng bằng dấu ngoặc và đặt dấu $+$'' trước mỗi ngoặc.
\item Cộng trừ các hệ số trong từng nhóm để thu gọn.
\end{enumerate}
Lưu ý: Đặc biệt cẩn thận khi trừ đa thức, luôn viết đa thức bị trừ trong ngoặc trước khi đổi dấu từng số hạng.

%% DANG 8D1?3-2
\textbf{Dạng 2. Tìm đa thức thỏa mãn đẳng thức cho trước}
\begin{enumerate}
\item Coi đa thức cần tìm $ như một số chưa biết và sử dụng quy tắc chuyển vế:
\begin{itemize}
\item Nếu  + A = B \implies \boxed{X = B - A}$.
\item Nếu  - A = B \implies \boxed{X = B + A}$.
\item Nếu  - X = B \implies \boxed{X = A - B}$.
\end{itemize}
\item Thực hiện phép tính cộng, trừ các đa thức ở vế phải rồi thu gọn kết quả.
\end{enumerate}

%% DANG 8D1?3-3
\textbf{Dạng 3. Bài toán thực tế dùng cộng, trừ đa thức}
\begin{enumerate}
\item Xác định các đại lượng hình học hoặc đời sống (chu vi, diện tích, chi phí, doanh thu\dots) theo các biến , y$.
\item Lập biểu thức đa thức biểu diễn từng đại lượng.
\item Thực hiện phép cộng hoặc trừ các đa thức để trả lời yêu cầu bài toán.
\item Thay giá trị số cụ thể của biến (nếu đề bài yêu cầu) để tính kết quả cuối cùng kèm đơn vị.
\end{enumerate}
